% ------------------------------------------------------------------------------
% Chapter 3
% ------------------------------------------------------------------------------
\chapter{Formal Statements} 
\label{cha:formal_statements}

Here we will introduce some useful logic followed by each statement, starting from the assumptions and then move on to the assertions.

\section{Logic}
\begin{description}
    \item[\texttt{past\_valid}]
        A helper register initialized to 0 that transitions to 1 after the first clock edge.
        It safely gates properties that evaluate historical data via the \texttt{\$past()} operator, preventing tool errors on cycle zero.
\end{description}

\section{Assumptions}
\begin{description}
    \item[\texttt{initial\_reset}] 
        Forces the \texttt{reset} signal to be active during the very first clock cycle (when \texttt{past\_valid} is 0) to ensure the FSM initializes cleanly.
    \item[\texttt{exclusive\_outcome}]
        Constrains the solver so that \texttt{chk\_ok} and \texttt{chk\_fail} are mutually exclusive and never asserted simultaneously.
        This could have been transformed into an assertion, but it would have tested the input instead of the dut itself.
\end{description}

\section{Assertions}
\begin{description}
    \item[\texttt{\_reset} \& \texttt{\_abort}]
        Verify that asserting a reset or an abort unconditionally forces the FSM back to the \texttt{IDLE} state on the next cycle with both output flags cleared.
    \item[\texttt{\_state}]
        Guarantees the FSM is mathematically restricted to the five valid encoded states, proving no illegal states can be reached.
    \item[State Progression] (\texttt{idle\_header}, \texttt{header\_payload}, \texttt{payload\_checksum}, \texttt{checkusm\_done}, \texttt{checksum\_idle}, \texttt{done\_idle}):
        Prove that the FSM correctly transitions to the next expected state when the appropriate progression flag is asserted.
    \item[State Hold] (\texttt{idle\_idle}, \texttt{header\_header}, \texttt{payload\_payload}, \texttt{checksum\_checksum}):
        Guarantee that the FSM retains its current state when no progression signal or abort is asserted.
    \item[Output Safety] (\texttt{safe\_valid\_pkt}, \texttt{safe\_error\_pkt}):
        Prove that the output flags are never asserted spuriously.
        They only pulse high if the FSM was actively in the \texttt{CHECKSUM} state during the previous cycle with the correct outcome flag.
    \item[Output Pulse] (\texttt{pulse\_valid\_pkt}, \texttt{pulse\_error\_pkt}):
        Verify that the output flags remain asserted for exactly one clock cycle, preventing it from a pulse behavior.
\end{description}